\documentclass[journal]{IEEEtran}

% Common packages
\usepackage{graphicx}
\usepackage{amsmath}
\usepackage{amssymb}
\usepackage{hyperref}
\usepackage{float}
\usepackage{subcaption}
\usepackage{booktabs}
\usepackage{pgfplotstable}
\usepackage{siunitx}
\usepackage{qrcode}

\pgfplotsset{compat=1.18}

\begin{document}

% ---------- TITLE & AUTHOR ----------
\title{Thermal Equation of State and Critical Point}
\author{Ibrahim H.I. ABUSHWISH \\
Istanbul University, Department of Physics \\
Instructor: \\
Course \& Section Number: PHYS3604}
\markboth{Physics Laboratory Reports, May 2026}{}

\maketitle

% ---------- ABSTRACT ----------
\begin{abstract}
    This report shows the derivation of th e Van der Waals constatns theoretically.
\end{abstract}

% ---------- INTRODUCTION ----------
\section{Introduction}
Provide physical background and motivation for the experiment. State the aim and what will be measured.

% ---------- THEORY ----------
\section{Theory}
The Van der Waals equation of state for a real gas is
\begin{equation}
P = \frac{RT}{V - b} - \frac{a}{V^2},
\label{eq:vdw}
\end{equation}
where $P$ is pressure, $T$ temperature, $V$ molar volume, $R$ the gas constant, and $a,b$ are constants that account for intermolecular attraction and finite molecular size respectively.

At the critical point $(P_c,T_c,V_c)$ the isotherm has a horizontal inflection, so
\begin{equation}
\left(\frac{\partial P}{\partial V}\right)_{T_c}=0,\qquad
\left(\frac{\partial^2 P}{\partial V^2}\right)_{T_c}=0.
\label{eq:critconds}
\end{equation}

Compute the derivatives from the Van der Waals form:
\begin{align}
\frac{\partial P}{\partial V} &= -\frac{RT}{(V-b)^2} + \frac{2a}{V^3},\\
\frac{\partial^2 P}{\partial V^2} &= \frac{2RT}{(V-b)^3} - \frac{6a}{V^4}.
\end{align}

Applying the critical conditions at $V=V_c, T=T_c$ gives:
\begin{align}
-\frac{RT_c}{(V_c-b)^2} + \frac{2a}{V_c^3} &= 0, \label{eq:crit1}\\
\frac{2RT_c}{(V_c-b)^3} - \frac{6a}{V_c^4} &= 0. \label{eq:crit2}
\end{align}

Dividing Eq.~\eqref{eq:crit2} by Eq.~\eqref{eq:crit1} yields
\begin{equation}
-\frac{2}{V_c-b} = \frac{3}{V_c},
\end{equation}
which gives
\begin{equation}
V_c = 3b.
\end{equation}

Substituting $V_c=3b$ into Eq.~\eqref{eq:crit1} leads to
\begin{equation}
T_c = \frac{8a}{27bR}.
\end{equation}
Evaluating the Van der Waals equation at the critical point produces
\begin{equation}
P_c = \frac{a}{27b^2}.
\end{equation}

Thus the standard critical-point relations are
\begin{equation}
V_c = 3b,\qquad P_c = \frac{a}{27b^2},\qquad T_c = \frac{8a}{27bR}.
\end{equation}

Solving for $a$ and $b$ from measurable critical constants $(P_c,T_c)$ gives
\begin{equation}
b = \frac{R T_c}{8P_c},\qquad a = \frac{27R^2 T_c^2}{64P_c}.
\end{equation}

The critical compressibility factor is
\begin{equation}
Z_c = \frac{P_c V_c}{R T_c} = \frac{3}{8}.
\end{equation}

These relations are used to determine Van der Waals constants from experimental critical data.
\end{document}